% concatenated from Minimal_Binary_Linear_Codes_of_Dimension_n+4_from_Partial_Spreads_and_Their_Dual_Access_Structures.tex
%Version 3.1 December 2024
% See section 11 of the User Manual for version history
%
%%%%%%%%%%%%%%%%%%%%%%%%%%%%%%%%%%%%%%%%%%%%%%%%%%%%%%%%%%%%%%%%%%%%%%
%%                                                                 %%
%% Please do not use \input{...} to include other tex files.       %%
%% Submit your LaTeX manuscript as one .tex document.              %%
%%                                                                 %%
%% All additional figures and files should be attached             %%
%% separately and not embedded in the \TeX\ document itself.       %%
%%                                                                 %%
%%%%%%%%%%%%%%%%%%%%%%%%%%%%%%%%%%%%%%%%%%%%%%%%%%%%%%%%%%%%%%%%%%%%%

%%\documentclass[referee,sn-basic]{sn-jnl}% referee option is meant for double line spacing

%%=======================================================%%
%% to print line numbers in the margin use lineno option %%
%%=======================================================%%

%%\documentclass[lineno,pdflatex,sn-basic]{sn-jnl}% Basic Springer Nature Reference Style/Chemistry Reference Style

%%=========================================================================================%%
%% the documentclass is set to pdflatex as default. You can delete it if not appropriate.  %%
%%=========================================================================================%%

%%\documentclass[sn-basic]{sn-jnl}% Basic Springer Nature Reference Style/Chemistry Reference Style

%%Note: the following reference styles support Namedate and Numbered referencing. By default the style follows the most common style. To switch between the options you can add or remove “Numbered” in the optional parenthesis. 
%%The option is available for: sn-basic.bst, sn-chicago.bst%  
 
%%\documentclass[pdflatex,sn-nature]{sn-jnl}% Style for submissions to Nature Portfolio journals
%%\documentclass[pdflatex,sn-basic]{sn-jnl}% Basic Springer Nature Reference Style/Chemistry Reference Style
\documentclass[pdflatex,sn-mathphys-num]{sn-jnl}% Math and Physical Sciences Numbered Reference Style
%%\documentclass[pdflatex,sn-mathphys-ay]{sn-jnl}% Math and Physical Sciences Author Year Reference Style
%%\documentclass[pdflatex,sn-aps]{sn-jnl}% American Physical Society (APS) Reference Style
%%\documentclass[pdflatex,sn-vancouver-num]{sn-jnl}% Vancouver Numbered Reference Style
%%\documentclass[pdflatex,sn-vancouver-ay]{sn-jnl}% Vancouver Author Year Reference Style
%%\documentclass[pdflatex,sn-apa]{sn-jnl}% APA Reference Style
%%\documentclass[pdflatex,sn-chicago]{sn-jnl}% Chicago-based Humanities Reference Style

%%%% Standard Packages
%%<additional latex packages if required can be included here>

\usepackage{graphicx}%
\usepackage{multirow}%
\usepackage{amsmath,amssymb,amsfonts}%
\usepackage{amsthm}%
\usepackage{mathrsfs}%
\usepackage[title]{appendix}%
\usepackage{xcolor}%
\usepackage{textcomp}%
\usepackage{manyfoot}%
\usepackage{booktabs}%
\usepackage{algorithm}%
\usepackage{algorithmicx}%
\usepackage{algpseudocode}%
\usepackage{listings}%
%%%%

%%%%%=============================================================================%%%%
%%%%  Remarks: This template is provided to aid authors with the preparation
%%%%  of original research articles intended for submission to journals published 
%%%%  by Springer Nature. The guidance has been prepared in partnership with 
%%%%  production teams to conform to Springer Nature technical requirements. 
%%%%  Editorial and presentation requirements differ among journal portfolios and 
%%%%  research disciplines. You may find sections in this template are irrelevant 
%%%%  to your work and are empowered to omit any such section if allowed by the 
%%%%  journal you intend to submit to. The submission guidelines and policies 
%%%%  of the journal take precedence. A detailed User Manual is available in the 
%%%%  template package for technical guidance.
%%%%%=============================================================================%%%%

%% as per the requirement new theorem styles can be included as shown below
\theoremstyle{thmstyleone}%
\newtheorem{theorem}{Theorem}%  meant for continuous numbers
%%\newtheorem{theorem}{Theorem}[section]% meant for sectionwise numbers
%% optional argument [theorem] produces theorem numbering sequence instead of independent numbers for Proposition
\newtheorem{proposition}[theorem]{Proposition}% 
\newtheorem{conjecture}{Conjecture}
\newtheorem{lemma}{Lemma}
%%\newtheorem{proposition}{Proposition}% to get separate numbers for theorem and proposition etc.

\theoremstyle{thmstyletwo}%
\newtheorem{example}{Example}%
\newtheorem{remark}{Remark}%

\theoremstyle{thmstylethree}%
\newtheorem{definition}{Definition}%
\newtheorem{corollary}{Corollary}


\raggedbottom
%%\unnumbered% uncomment this for unnumbered level heads

\begin{document}

\title[Minimal Binary Linear Codes of Dimension $n+4$ from Partial Spreads and Their Dual Access Structures]{Minimal Binary Linear Codes of Dimension $n+4$ from Partial Spreads and Their Dual Access Structures}

%%=============================================================%%
%% GivenName	-> \fnm{Joergen W.}
%% Particle	-> \spfx{van der} -> surname prefix
%% FamilyName	-> \sur{Ploeg}
%% Suffix	-> \sfx{IV}
%% \author*[1,2]{\fnm{Joergen W.} \spfx{van der} \sur{Ploeg} 
%%  \sfx{IV}}\email{iauthor@gmail.com}
%%=============================================================%%

\author[1]{\fnm{Apurba} \sur{Sarkar}}\email{apurbasarkar065@gmail.com}

\author[1]{\fnm{Kalyan} \sur{Hansda}}\email{kalyanh4@gmail.com}
\equalcont{These authors contributed equally to this work.}

\author*[2]{\fnm{Makhan} \sur{Maji}}\email{makhan2maths@gmail.com}
\equalcont{These authors contributed equally to this work.}

\affil[1]{\orgdiv{Department of Mathematics}, \orgname{Visva-Bharati}, \orgaddress{\city{Santiniketen}, \postcode{731235}, \state{West Bengal}, \country{India}}}

\affil*[2]{\orgname{Indian  Institute of Technology Madras}, \orgaddress{\city{Chennai}, \postcode{600036}, \state{TN}, \country{India}}}

%%==================================%%
%% Sample for unstructured abstract %%
%%==================================%%

\abstract{Minimal linear codes have significant applications in secret sharing schemes, secure multi-party computation, and cryptography. In this paper, we propose a generic construction of a new family of minimal binary linear codes with dimension $n+4$ from a special class of Boolean functions. By leveraging the geometric properties of partial spreads in finite fields, we determine the explicit weight distribution and weight enumerator of the constructed codes. Furthermore, we derive a necessary and sufficient condition for these codes to be minimal, and establish that the proposed family yields minimal codes that structurally violate the well-known Ashikhmin-Barg condition, making them highly desirable for advanced communication systems.}


\keywords{Module, Linear code, Support, Minimal linear code}

%%\pacs[JEL Classification]{D8, H51}

\pacs[MSC Classification]{94B05; 11T71}

\maketitle

\section{Introduction and Preliminaries}

In modern communication architectures and multi-party cryptographic frameworks, error-correcting linear codes serve a dual purpose: ensuring robust data reliability over noisy channels and providing foundational primitives for cryptographic protocols \cite{LingXing2004}. Let $\mathbb{F}_{2}^n$ denote the $n$-dimensional vector space over the binary field $\mathbb{F}_{2}$. An $[N, K, D]$ binary linear code $\mathcal{M}$ is a $K$-dimensional subspace of $\mathbb{F}_{2}^N$ with minimum Hamming distance $D$ \cite{LingXing2004}. For a codeword $m = (m_1, m_2, \dots, m_N) \in \mathcal{M}$, its coordinate support is $\text{Supp}(m) = \{1 \le l \le N \mid m_l \neq 0\}$, and its Hamming weight is $wt(m) = |\text{Supp}(m)|$.

A crucial subclass of linear codes is the family of \emph{minimal linear codes}. A non-zero codeword $m \in \mathcal{M}$ is called minimal if its support does not properly contain the support of any other linearly independent codeword in $\mathcal{M}$; that is, $\text{Supp}(m') \subseteq \text{Supp}(m)$ implies $m' = m$ over $\mathbb{F}_2$. Minimal linear codes form the backbone of Massey's secret sharing schemes (SSS) \cite{Massey1993}, where the minimal codewords of the dual code $\mathcal{M}^\perp$ uniquely dictate the minimal access structures.

Historically, verifying code minimality relied on the Ashikhmin--Barg (AB) sufficient condition \cite{AshikhminBarg1998}, which guarantees minimality if the weight ratio satisfies:
\begin{equation}\label{eq:ash_barg}
\frac{w_{\min}}{w_{\max}} > \frac{1}{2}.
\end{equation}
While simple, the AB condition tightly clusters codeword weights, producing rigid secret sharing access structures with restricted operational flexibility. This limitation motivated a shift toward constructing \emph{non-Ashikhmin--Barg minimal linear codes}—codes satisfying $w_{\min}/w_{\max} \le 1/2$ that remain minimal due to deep algebraic and geometric properties \cite{DingHengZhou2018, HengDingZhou2018, LuWuCao2021}. Significant progress has been made using simplicial complexes, special functions, and Boolean function combinations \cite{CarletDingYuan2005, Ding2016, LiYue2020, MesnagerQianCaoYuan2023, YuanDing2005}.

Recently, Liu and Liao \cite{LiuLiao2022} developed a generic framework using Boolean functions supported on partial spreads. However, existing constructions remain primarily confined to lower-dimensional functional extensions ($K = n+1, n+2, n+3$), leaving higher-dimensional regimes unexplored. In this paper, we construct a new class of minimal binary linear codes featuring an expanded dimension of $K = n + 4$ over block length $N = 2^n - 1$ using a composite space of four partial spread indicator functions $\mathcal{F} = \text{span}\{\psi_1, \psi_2, \psi_3, \psi_4\}$.

Rather than a routine parameter expansion, this work offers fundamental algebraic and cryptographic contributions to the literature:

\begin{enumerate}
    \item \textbf{Circumvention of the Ashikhmin--Barg Ceiling at Dimension $n+4$:} Expanding the functional basis to four blocks significantly increases $w_{\max}$, pushing $w_{\min}/w_{\max}$ strictly below $1/2$. We prove that our combinatorial non-containment design rules (Conditions C1--C3) act as an exact geometric sieve, ensuring code minimality at dimension $n+4$ independently of global weight bounds.

    \item \textbf{Resolution of 15-Pair Combinatorial Saturation over $\mathbb{F}_2^4$:} A four-function basis generates $2^4 - 1 = 15$ non-zero Boolean combinations $\phi \in \mathcal{F}$. Proving minimality requires showing non-containment across all $15 \times 15 = 225$ function pair interactions simultaneously, establishing $n + 4$ as a natural combinatorial boundary for partial spread constructions.

    \item \textbf{Resolution of the Ashikhmin--Barg Rigidity Bottleneck in SSS:} While classical AB-compliant minimal codes yield rigid access structures with tightly clustered set cardinalities, our non-AB $n+4$ codes induce a multi-threshold access spectrum with a broad operational span $\Delta \ge 2^{n-1} + 2^{m-2} - 1$, enabling flexible, multi-tiered authorization hierarchies.

    \item \textbf{Quadrupled Authorization Density ($2^{n+3}$):} Over block length $N = 2^n - 1$, our $K = n+4$ construction expands the total number of minimal access sets to $|\Gamma| = 2^{n+3}$, quadrupling the authorization pathways available in $(n+2)$-dimensional schemes \cite{LiuLiao2022} and significantly enhancing network fault tolerance against participant dropouts.

    \item \textbf{Statistical Variance, Overlap Bounds, and Throughput Gain:} We prove that $Var(|\mathcal{P}_A|) > 0$ across eight distinct set-size tiers and derive a tight upper bound on participant overlap ($|\mathcal{P}_{A_1} \cap \mathcal{P}_{A_2}|$), guaranteeing formal immunity against statistical share approximation and targeted collusion attacks. Furthermore, our scheme achieves a $20\%$ relative information rate throughput enhancement over prior $(n+2)$-dimensional constructions.
\end{enumerate}

The remainder of this paper is organized as follows. Section 2 presents the generic construction of $\mathcal{M}_{(\psi_a)}$ and derives Walsh--Hadamard minimality criteria. Section 3 details the partial spread construction, establishes exact weight distributions, presents minimality proofs for $n \ge 8$, and provides concrete numerical examples. Section 4 explores applications in secret sharing schemes, establishing access set bounds, throughput gains, coalition immunity, and overlap limits.



\section{Minimal binary linear codes with dimension $n + 4$ from a class of Boolean functions}
In this section, we present a generic construction of binary linear codes with a dimension of $n + 4$ using a specialized class of Boolean functions. Furthermore, we establish a necessary and sufficient condition based on Walsh-Hadamard transform coefficients for the constructed codes to achieve minimality.

Let $\psi_a$ for $a = 1,2,3,4$ be Boolean functions defined on $\mathbb{F}_2^n$. We denote by $\mathcal{F}$ the linear span of these functions over $\mathbb{F}_2$:
\begin{equation}\label{eq:F_space}
\mathcal{F} = \left\{\sum_{k=1}^{4} \alpha_k \psi_k \;\middle|\; \alpha_k \in \mathbb{F}_2 \right\},
\end{equation}
such that the following fundamental conditions hold simultaneously:
\begin{enumerate}
    \item $\phi$ is a non-zero Boolean function for all $\phi \in \mathcal{F}$.
    \item $\phi(\mathbf{0}) = 0$ for all $\phi \in \mathcal{F}$.
    \item $\phi_1 \neq \phi_2$ for all distinct $\phi_1, \phi_2 \in \mathcal{F}$.
\end{enumerate}

We construct a linear code $\mathcal{M}_{(\psi_a)} = \mathcal{M}(\psi_1, \psi_2, \psi_3, \psi_4)$ of length $2^n-1$ as an extension of the classical simplex code. The codewords are defined by evaluating the function combinations over all non-zero vectors:
\begin{equation}\label{code}
\mathcal{M}_{(\psi_a)} = \left\{ m_{(\psi_a)}(\mathbf{\beta}) = \left( \sum_{k=1}^{4} \alpha_k \psi_k(\mathbf{x}) + \mathbf{\beta} \cdot \mathbf{x} \right)_{\mathbf{x} \in \mathbb{F}_2^{n*}} \;\middle|\; \alpha_k \in \mathbb{F}_2,\; \mathbf{\beta} \in \mathbb{F}_2^n \right\}.
\end{equation}

\begin{theorem}\label{C_weight}
Let $\mathcal{M}_{(\psi_a)}$ be the binary code defined in \eqref{code}. If $\phi(\mathbf{x})\neq \mathbf{v}\cdot \mathbf{x}$ for all $\mathbf{v} \in \mathbb{F}_2^n$ and all $\phi \in \mathcal{F}$, then $\mathcal{M}_{(\psi_a)}$ is a $[2^n-1, n+4]$ binary linear code whose weight distribution is given by the multiset union:
\begin{equation}\label{eq:multiset_weight}
\Omega = \bigcup_{\phi \in \mathcal{F}} \left\{ \frac{2^{n} - \hat{\phi}(\mathbf{\beta})}{2} \;\middle|\; \mathbf{\beta} \in \mathbb{F}_{2}^{n} \right\} \;\cup\; \{2^{n-1} \mid \mathbf{\beta} \in \mathbb{F}_{2}^{n*}\} \;\cup\; \{0\}.
\end{equation}
\end{theorem}

\begin{proof}
The length of $\mathcal{M}_{(\psi_a)}$ is directly determined by the cardinality of $\mathbb{F}_2^{n*}$, which is $2^n - 1$. Let $\{\mathbf{u}_j \mid j = 1,2,\ldots,n\}$ be a basis of the vector space $\mathbb{F}_2^n$. To establish the dimension, we show that the set $\mathcal{B} = \{m_{(\psi_a)}(\mathbf{u}_j) \mid 1 \leq j \leq n\} \cup \{m_{\psi_a}(\mathbf{0}) \mid 1 \leq a \leq 4\}$ forms a basis for $\mathcal{M}_{(\psi_a)}$. Setting a linear combination to zero gives:
\begin{align*}
&\sum_{j=1}^{n} \delta_j m_{(\psi_a)}(\mathbf{u}_j) + \sum_{k=1}^{4} \delta_{n+k} m_{\psi_k}(\mathbf{0}) = \mathbf{0}\\
\Rightarrow & \sum_{j=1}^{n} \delta_j 
\left(
\sum_{k=1}^{4} \psi_k(\mathbf{x}) + \mathbf{u}_j \cdot \mathbf{x}
\right)_{\mathbf{x} \in \mathbb{F}_2^{n*}}
+ \sum_{k=1}^{4} \delta_{n+k} (\psi_k(\mathbf{x}))_{\mathbf{x} \in \mathbb{F}_2^{n*}} = \mathbf{0}\\
\Rightarrow &\sum_{j=1}^{n} \delta_j (\mathbf{u}_j \cdot \mathbf{x}) + \sum_{k=1}^{4} \gamma_k \psi_k(\mathbf{x}) = 0, 
\end{align*}
where $\gamma_k = \sum_{j=1}^{n} \delta_j + \delta_{n+k}$. If $\gamma_k \neq 0$ for any $k$, it implies that some linear combination $\phi \in \mathcal{F}$ equals an inner product $\mathbf{v} \cdot \mathbf{x}$, which contradicts our initial assumption. Thus, $\gamma_k = 0$ for all $k$, which directly forces $\delta_j = 0$ for all $1 \leq j \leq n+4$. Since no subset of size $n+5$ can be linearly independent within this construction, the dimension of the code is exactly $n+4$. 

To establish the weight distribution, we evaluate the Hamming weight of the codewords using the Walsh-Hadamard transform. For the case where $\alpha_k = 1$ for all $k = 1,2,3,4$, the codeword corresponds to the evaluation vector of the full functional sum $\sum_{k=1}^{4} \psi_k$. Its Walsh-Hadamard transform at a shift vector $\mathbf{\beta} \in \mathbb{F}_2^n$ is given by:
\begin{align*}
\widehat{\sum_{k=1}^{4} \psi_k}(\mathbf{\beta}) 
&= \sum_{\mathbf{x} \in \mathbb{F}_2^n} (-1)^{\sum_{k=1}^{4} \psi_k(\mathbf{x}) + \mathbf{\beta} \cdot \mathbf{x}} \\
&= \left(2^n - \mathrm{wt}(m_{(\psi_a)}(\mathbf{\beta}))\right) - \mathrm{wt}(m_{(\psi_a)}(\mathbf{\beta})) \\
&= 2^n - 2\,\mathrm{wt}(m_{(\psi_a)}(\mathbf{\beta})).
\end{align*}
Isolating the weight term yields the explicit coordinate weight expression:
\[
\mathrm{wt}(m_{(\psi_a)}(\mathbf{\beta})) = \frac{2^n - \widehat{\sum_{k=1}^{4} \psi_k}(\mathbf{\beta})}{2}.
\]
The weight enumerations for all remaining configurations of the coefficients $\alpha_k \in \mathbb{F}_2$ are derived in an exactly analogous manner.
\end{proof}

We now establish a necessary and sufficient condition for the linear code $\mathcal{M}_{(\psi_a)}$ to achieve minimality based on its Walsh-Hadamard transform coefficients.
\begin{theorem}\label{Thm_Minimal_Cond}
The binary linear code $\mathcal{M}_{(\psi_a)}$ defined in \eqref{code} is minimal if and only if the following two conditions hold simultaneously for all vectors $\mathbf{x},\mathbf{y} \in \mathbb{F}_2^n$:
\begin{enumerate}
    \item For any $\phi_1,\phi_2 \in \mathcal{F}$ and $\mathbf{x} \neq \mathbf{y}$:
    \[ \widehat{\phi_1}(\mathbf{x})+\widehat{\phi_2}(\mathbf{y})\neq 2^n \quad \text{and} \quad \widehat{\phi_1}(\mathbf{x})-\widehat{\phi_2}(\mathbf{y})\neq 2^n. \tag{2.3} \]
    \item For any distinct $\phi_1,\phi_2 \in \mathcal{F}$:
    \[ \widehat{\phi_1}(\mathbf{x}+\mathbf{y})+\widehat{\phi_2}(\mathbf{x})-\widehat{\phi_1+\phi_2}(\mathbf{y})\neq 2^n. \tag{2.4} \]
\end{enumerate}
\end{theorem}

\begin{proof}
Every codeword $\mathbf{m}_{(\psi_a)}(\mathbf{\beta}) \in \mathcal{M}_{(\psi_a)}$ can be uniquely decomposed as follows:
\[
m_{(\psi_a)}(\mathbf{\beta})=\sum_{k=1}^{4}\alpha_k \mathbf{\psi}_k + c(\mathbf{\beta}),
\]
where $\mathbf{\psi}_a = \bigl(\psi_a(\mathbf{x})\bigr)_{\mathbf{x} \in \mathbb{F}_2^{n*}}$ and $c(\mathbf{\beta})$ is an element of the subcode:
\[
\mathcal{C}=\left\{c(\mathbf{\beta})=\bigl(\mathbf{\beta}\cdot \mathbf{x}\bigr)_{\mathbf{x} \in \mathbb{F}_2^{n*}} \mid \mathbf{\beta}\in \mathbb{F}_2^n\right\}.
\]
Here, $\mathcal{C}$ is a simplex code with parameters $[2^n - 1, n, 2^{n-1}]$. Thus, for any non-zero linear shift vector, $\mathrm{wt}(c(\mathbf{\beta})) = 2^{n-1}$. 

We verify that no improper support containment occurs across all structural subcases:

\textbf{Case I:} Let both codewords share an identical linear shift vector component such that $\mathbf{m}_1, \mathbf{m}_2 \in \left\{\sum_{k=1}^{4} \alpha_k \mathbf{\psi}_k + c(\mathbf{\beta}) \mid \alpha_k \in \mathbb{F}_2 \right\}$. 

\textbf{(a)} Let $\mathbf{m}_1 = c(\mathbf{\beta})$ and $\mathbf{m}_2 = \sum_{k=1}^{4} \mathbf{\psi}_k + c(\mathbf{\beta})$. Utilizing the weight distributions established via Theorem~\ref{C_weight}, the coordinate ordering evaluates to:
\begin{align*}
\mathbf{m}_1 \preceq \mathbf{m}_2 &\Leftrightarrow 
\mathrm{wt}\!\left(\sum_{k=1}^{4} \mathbf{\psi}_k\right) 
= \mathrm{wt}\!\left(\sum_{k=1}^{4} \mathbf{\psi}_k + c(\mathbf{\beta})\right) - 2^{n-1}\\
& \Leftrightarrow \widehat{\sum_{k=1}^{4} \psi_k}(\mathbf{0}) - \widehat{\sum_{k=1}^{4} \psi_k}(\mathbf{\beta}) = 2^n.
\end{align*}
Testing the reverse ordering configuration yields:
\[
\mathbf{m}_2 \preceq \mathbf{m}_1 \Leftrightarrow 
\mathrm{wt}\!\left(\sum_{k=1}^{4} \mathbf{\psi}_k\right) 
= 2^{n-1} - \mathrm{wt}\!\left(\sum_{k=1}^{4} \mathbf{\psi}_k + c(\mathbf{\beta})\right)
\Leftrightarrow \widehat{\sum_{k=1}^{4} \psi_k}(\mathbf{0}) + \widehat{\sum_{k=1}^{4} \psi_k}(\mathbf{\beta}) = 2^n.
\]

\textbf{(b)} Let $\mathbf{m}_1 = \mathbf{\psi}_a + \mathbf{\psi}_b + c(\mathbf{\beta})$ and $\mathbf{m}_2 = \mathbf{\psi}_c + \mathbf{\psi}_d + c(\mathbf{\beta})$ for $1 \le a,b,c,d \le 4$ with $a \neq b$ and $c \neq d$. The coordinate containments evaluate to:
\begin{align*}
&\mathbf{m}_1 \preceq \mathbf{m}_2 \Leftrightarrow 
\mathrm{wt}(\mathbf{\psi}_a + \mathbf{\psi}_b + \mathbf{\psi}_c + \mathbf{\psi}_d) \\
&= \mathrm{wt}(\mathbf{\psi}_c + \mathbf{\psi}_d + c(\mathbf{\beta})) - \mathrm{wt}(\mathbf{\psi}_a + \mathbf{\psi}_b + c(\mathbf{\beta}))\\
&\Leftrightarrow \widehat{\psi_a + \psi_b + \psi_c + \psi_d}(\mathbf{0}) + \widehat{\psi_a + \psi_b}(\mathbf{\beta}) - \widehat{\psi_c + \psi_d}(\mathbf{\beta}) = 2^n.
\end{align*}
Symmetrically, checking the opposite directional relation yields:
\[
\mathbf{m}_2 \preceq \mathbf{m}_1 \Leftrightarrow 
\mathrm{wt}(\mathbf{\psi}_a + \mathbf{\psi}_b + \mathbf{\psi}_c + \mathbf{\psi}_d) 
= \mathrm{wt}(\mathbf{\psi}_a + \mathbf{\psi}_b + c(\mathbf{\beta})) - \mathrm{wt}(\mathbf{\psi}_c + \mathbf{\psi}_d + c(\mathbf{\beta}))
\]
\[
\Leftrightarrow \widehat{\psi_a + \psi_b + \psi_c + \psi_d}(\mathbf{0}) + \widehat{\psi_c + \psi_d}(\mathbf{\beta}) - \widehat{\psi_a + \psi_b}(\mathbf{\beta}) = 2^n.
\]

\textbf{Case II:} Let the codewords be defined over distinct linear shift vectors such that $\mathbf{\beta}_1 \neq \mathbf{\beta}_2$:
\[
\mathbf{m}_1 \in \left\{\sum_{k=1}^{4} \alpha_k \mathbf{\psi}_k + c(\mathbf{\beta}_1) \;\middle|\; \alpha_k \in \mathbb{F}_2 \right\} \quad \text{and} \quad \mathbf{m}_2 \in \left\{\sum_{k=1}^{4} \alpha_k \mathbf{\psi}_k + c(\mathbf{\beta}_2) \;\middle|\; \alpha_k \in \mathbb{F}_2 \right\}.
\]

\textbf{(a)} Let $\mathbf{m}_1 = c(\mathbf{\beta}_1)$ and $\mathbf{m}_2 = \sum_{k=1}^{4} \mathbf{\psi}_k + c(\mathbf{\beta}_2)$. Evaluating the translations yields:
\[
\mathbf{m}_1 \preceq \mathbf{m}_2 \Leftrightarrow 
\mathrm{wt}\!\left(\sum_{k=1}^{4} \mathbf{\psi}_k + c(\mathbf{\beta}_1 + \mathbf{\beta}_2)\right)
= \mathrm{wt}\!\left(\sum_{k=1}^{4} \mathbf{\psi}_k + c(\mathbf{\beta}_2)\right) - 2^{n-1}
\]
\[
\Leftrightarrow \widehat{\sum_{k=1}^{4} \psi_k}(\mathbf{\beta}_1 + \mathbf{\beta}_2) - \widehat{\sum_{k=1}^{4} \psi_k}(\mathbf{\beta}_2) = 2^n.
\]
Reversing the containment evaluation gives:
\[
\mathbf{m}_2 \preceq \mathbf{m}_1 \Leftrightarrow 
\mathrm{wt}\!\left(\sum_{k=1}^{4} \mathbf{\psi}_k + c(\mathbf{\beta}_1 + \mathbf{\beta}_2)\right)
= 2^{n-1} - \mathrm{wt}\!\left(\sum_{k=1}^{4} \mathbf{\psi}_k + c(\mathbf{\beta}_2)\right)
\]
\[
\Leftrightarrow \widehat{\sum_{k=1}^{4} \psi_k}(\mathbf{\beta}_1 + \mathbf{\beta}_2) + \widehat{\sum_{k=1}^{4} \psi_k}(\mathbf{\beta}_2) = 2^n.
\]

\textbf{(b)} Let $\mathbf{m}_1 = \boldsymbol{\phi} + c(\mathbf{\beta}_1)$ and $\mathbf{m}_2 = \boldsymbol{\phi} + c(\mathbf{\beta}_2)$ with $\boldsymbol{\phi} = \mathbf{\psi}_a + \mathbf{\psi}_b + \mathbf{\psi}_c$ for $1\leq a,b,c \leq 4$ and $a \neq b \neq c$. Evaluating the vector shifts maps to:
\[
\mathbf{m}_1 \preceq \mathbf{m}_2 \Leftrightarrow 
\mathrm{wt}(c(\mathbf{\beta}_1 + \mathbf{\beta}_2)) 
= \mathrm{wt}(\mathbf{\psi}_a + \mathbf{\psi}_b + \mathbf{\psi}_c + c(\mathbf{\beta}_2)) - \mathrm{wt}(\mathbf{\psi}_a + \mathbf{\psi}_b + \mathbf{\psi}_c + c(\mathbf{\beta}_1))
\]
\[
\Leftrightarrow \widehat{\psi_a + \psi_b + \psi_c}(\mathbf{\beta}_1) - \widehat{\psi_a + \psi_b + \psi_c}(\mathbf{\beta}_2) = 2^n.
\]
Evaluating the opposite structural direction yields:
\[
\mathbf{m}_2 \preceq \mathbf{m}_1 \Leftrightarrow 
\mathrm{wt}(c(\mathbf{\beta}_1 + \mathbf{\beta}_2)) 
= \mathrm{wt}(\mathbf{\psi}_a + \mathbf{\psi}_b + \mathbf{\psi}_c+ c(\mathbf{\beta}_1)) - \mathrm{wt}(\mathbf{\psi}_a + \mathbf{\psi}_b + \mathbf{\psi}_c + c(\mathbf{\beta}_2))
\]
\[
\Leftrightarrow \widehat{\psi_a + \psi_b + \psi_c}(\mathbf{\beta}_2) - \widehat{\psi_a + \psi_b + \psi_c}(\mathbf{\beta}_1) = 2^n.
\]

\textbf{(c)} Let $\mathbf{m}_1 = \boldsymbol{\phi}_1 + c(\mathbf{\beta}_1)$ and $\mathbf{m}_2 = \boldsymbol{\phi}_2 + c(\mathbf{\beta}_2)$ with $\boldsymbol{\phi}_1 = \mathbf{\psi}_a + \mathbf{\psi}_b$ and $\boldsymbol{\phi}_2 = \mathbf{\psi}_c + \mathbf{\psi}_d$ for distinct $\boldsymbol{\phi}_1 \neq \boldsymbol{\phi}_2$. Resolving the underlying code constraints gives:
\[
\mathbf{m}_1 \preceq \mathbf{m}_2 \Leftrightarrow 
\mathrm{wt}\!\left(\sum_{k=1}^{4} \mathbf{\psi}_k + c(\mathbf{\beta}_1 + \mathbf{\beta}_2)\right) 
= \mathrm{wt}(\mathbf{\psi}_c + \mathbf{\psi}_d + c(\mathbf{\beta}_2)) - \mathrm{wt}(\mathbf{\psi}_a + \mathbf{\psi}_b + c(\mathbf{\beta}_1))
\]
\[
\Leftrightarrow \widehat{\sum_{k=1}^{4} \psi_k}(\mathbf{\beta}_1 + \mathbf{\beta}_2) + \widehat{\psi_a + \psi_b}(\mathbf{\beta}_1) - \widehat{\psi_c + \psi_d}(\mathbf{\beta}_2) = 2^n,
\]
and symmetrically for the inverse containment boundary:
\[
\mathbf{m}_2 \preceq \mathbf{m}_1 \Leftrightarrow 
\mathrm{wt}\!\left(\sum_{k=1}^{4} \mathbf{\psi}_k + c(\mathbf{\beta}_1 + \mathbf{\beta}_2)\right) 
= \mathrm{wt}(\mathbf{\psi}_a + \mathbf{\psi}_b + c(\mathbf{\beta}_1)) - \mathrm{wt}(\mathbf{\psi}_c + \mathbf{\psi}_d + c(\mathbf{\beta}_2))
\]
\[
\Leftrightarrow\widehat{\sum_{k=1}^{4} \psi_k}(\mathbf{\beta}_1 + \mathbf{\beta}_2) +   \widehat{\psi_c + \psi_d}(\mathbf{\beta}_2) - \widehat{\psi_a + \psi_b}(\mathbf{\beta}_1) = 2^n.
\]
Because none of these improper algebraic equalities can occur under the simultaneous constraints of equations (2.3) and (2.4), every non-zero codeword is minimal.
\end{proof}




\section{A family of minimal binary linear codes with dimension $n+4$}

 Let $n$ be a positive even integer and let $m=\frac{n}{2}$. A partial spread of order $\kappa$ in $\mathbb{F}_2^n$ is a collection of $m$-dimensional subspaces $\{V_1, V_2, \ldots, V_\kappa\}$ such that $V_a \cap V_b = \{\mathbf{0}\}$ for all $a \neq b$. Note that, $\kappa\leq 2^{t}+1$.

For each $1 \leq \ell \leq \kappa$, we define an Boolean function $\eta_\ell$ on $\mathbb{F}_2^n$ as:
\begin{equation}\label{eq:g_def}
\eta_\ell(\mathbf{x})=
\begin{cases}
1, & \text{if } \mathbf{x} \in V_\ell \setminus \{\mathbf{0}\},\\
0, & \text{otherwise.}
\end{cases} \tag{3.1}
\end{equation}
Let $\mathcal{S}_1, \mathcal{S}_2, \mathcal{S}_3, \mathcal{S}_4$ be distinct, non-empty subsets of $\{1,2,\ldots,\kappa\}$, with respective cardinalities $|\mathcal{S}_a| = s_a$. We define our boolean functions $\psi_a$ ($1 \leq a \leq 4$) as:
\begin{equation}\label{eq:f_def}
\psi_a = \sum_{\ell \in \mathcal{S}_a} \eta_\ell. \tag{3.2}
\end{equation}
It follows from the disjoint nature of the partial spread elements that any linear combination $\phi \in \mathcal{F}$ can be represented cleanly via symmetric differences of the index sets:
\begin{equation}\label{eq:f_comb}
\psi_a+\psi_b=\sum_{\ell\in \mathcal{S}_a\triangle \mathcal{S}_b}\eta_\ell, \quad \psi_a+\psi_b+\psi_c=\sum_{\ell\in \mathcal{S}_a\triangle \mathcal{S}_b\triangle \mathcal{S}_c}\eta_\ell, \quad \text{and} \quad \sum_{k=1}^{4}\psi_k=\sum_{\ell\in \mathcal{S}_1\triangle \mathcal{S}_2\triangle \mathcal{S}_3\triangle \mathcal{S}_4}\eta_\ell. \tag{3.3}
\end{equation}
We denote the cardinalities of these symmetric differences by $\sigma_{ab} = |\mathcal{S}_a \triangle \mathcal{S}_b|$, $\sigma_{abc} = |\mathcal{S}_a \triangle \mathcal{S}_b \triangle \mathcal{S}_c|$, and $\sigma_{1234} = |\mathcal{S}_1 \triangle \mathcal{S}_2 \triangle \mathcal{S}_3 \triangle \mathcal{S}_4|$.

\begin{theorem}\label{Thm_Spread_Weight}
Let the  sets $\mathcal{S}_a$ be distinct and satisfy $\mathcal{S}_1 \triangle \mathcal{S}_2 \triangle \mathcal{S}_3 \triangle \mathcal{S}_4 \neq \emptyset$. Then, the code $\mathcal{M}_{(\psi_a)}$ generated by the functions in \eqref{eq:f_def} is a $[2^n-1, n+4]$ binary linear code whose explicit weight parameters and multiplicities are uniquely determined by Table 1.
\end{theorem}

\begin{proof}
Let $\mathbf{\beta} \in \mathbb{F}_2^n$. The subspaces $V_\ell$ for $1 \leq \ell \leq \kappa$ are pairwise disjoint and $\dim(V_\ell) = m$. For any fixed index $\ell$, the sub-sum character evaluates to:
\[
\sum_{\mathbf{x} \in V_\ell} (-1)^{\mathbf{\beta} \cdot \mathbf{x}} = 
\begin{cases} 
2^m, & \text{if } \mathbf{\beta} \in V_\ell^\perp, \\ 
0, & \text{if } \mathbf{\beta} \notin V_\ell^\perp. 
\end{cases}
\]
Since $\psi_a(\mathbf{x}) = \sum_{\ell \in \mathcal{S}_a} \eta_\ell(\mathbf{x})$, where $\eta_\ell(\mathbf{x}) = 1$ if $\mathbf{x} \in V_\ell \setminus \{\mathbf{0}\}$ and $0$ otherwise, we have:
\begin{align*}
\widehat{\psi_a}(\mathbf{\beta}) &= \sum_{\mathbf{x} \in \mathbb{F}_2^n} (-1)^{\psi_a(\mathbf{x}) + \mathbf{\beta} \cdot \mathbf{x}} \\
&= \sum_{\mathbf{x} \in \mathbb{F}_2^n \setminus \bigcup_{\ell \in \mathcal{S}_a} V_\ell} (-1)^{\mathbf{\beta} \cdot \mathbf{x}} + \sum_{\ell \in \mathcal{S}_a} \sum_{\mathbf{x} \in V_\ell \setminus \{\mathbf{0}\}} (-1)^{1 + \mathbf{\beta} \cdot \mathbf{x}} + \sum_{\ell\in\mathcal{S}_a}(-1)^{\beta\cdot\mathbf{0}}\\
&= \sum_{\mathbf{x} \in \mathbb{F}_2^n} (-1)^{\mathbf{\beta} \cdot \mathbf{x}} - \sum_{\ell \in \mathcal{S}_a} \sum_{\mathbf{x} \in V_\ell} (-1)^{\mathbf{\beta} \cdot \mathbf{x}} - \sum_{\ell \in \mathcal{S}_a} \left[\sum_{\mathbf{x} \in V_\ell} (-1)^{\mathbf{\beta} \cdot \mathbf{x}} -(-1)^{\mathbf{\beta} \cdot \mathbf{0}} \right]+\sum_{\ell\in\mathcal{S}_a}(-1)^{\beta\cdot\mathbf{0}} \\
&= \sum_{\mathbf{x} \in \mathbb{F}_2^n} (-1)^{\mathbf{\beta} \cdot \mathbf{x}} - 2 \sum_{\ell \in \mathcal{S}_a} \sum_{\mathbf{x} \in V_\ell} (-1)^{\mathbf{\beta} \cdot \mathbf{x}} + 2\sum_{\ell \in \mathcal{S}_a} 1\\
&= \sum_{\mathbf{x} \in \mathbb{F}_2^n} (-1)^{\mathbf{\beta} \cdot \mathbf{x}} - 2 \sum_{\ell \in \mathcal{S}_a} \sum_{\mathbf{x} \in V_\ell} (-1)^{\mathbf{\beta} \cdot \mathbf{x}} + 2s_a.
\end{align*}

We establish the values of $\widehat{\psi_a}(\mathbf{\beta})$ by analyzing the orthogonal spaces across three cases:

\textbf{Case 1:} $\mathbf{\beta} = \mathbf{0}$.
\begin{align*}
\widehat{\psi_a}(\mathbf{0}) &= 2^n - 2 \sum_{\ell \in \mathcal{S}_a} 2^m + 2s_a \\
&= 2^n - 2s_a 2^m + 2s_a \\
&= 2^n - 2s_a(2^m - 1).
\end{align*}

\textbf{Case 2:} $\mathbf{\beta} \neq \mathbf{0}$ and $\mathbf{\beta} \notin V_\ell^\perp$ for all $\ell \in \mathcal{S}_a$.
Since $\mathbf{\beta} \neq \mathbf{0}$, t$\sum_{\mathbf{x} \in \mathbb{F}_2^n} (-1)^{\mathbf{\beta} \cdot \mathbf{x}} = 0$, and the subspace character sums vanish for all $\ell \in \mathcal{S}_a$. Thus:
\begin{align*}
\widehat{\psi_a}(\mathbf{\beta}) &= 0 - 2 \sum_{\ell \in \mathcal{S}_a} 0 + 2\sum_{\ell \in \mathcal{S}_a} 1 \\
&= 2s_a.
\end{align*}

\textbf{Case 3:} $\mathbf{\beta} \neq \mathbf{0}$ and $\mathbf{\beta} \in V_{\ell_0}^\perp$ for a unique index $\ell_0 \in \mathcal{S}_a$. 
Since $V_\ell \cap V_k = \{\mathbf{0}\}$ for $\ell \neq k$, it follows that $\dim(V_\ell + V_k) = 2m = n$, which guarantees $V_\ell^\perp \cap V_k^\perp = \{\mathbf{0}\}$. Thus, $\mathbf{\beta}$ can belong to at most one $V_{\ell_0}^\perp \setminus \{\mathbf{0}\}$.
\begin{align*}
\widehat{\psi_a}(\mathbf{\beta}) &= 0 - 2 \left( \sum_{\ell \in \mathcal{S}_a \setminus \{\ell_0\}} 0 + 2^m \right) + 2s_a \\
&= -2^{m+1} + 2s_a.
\end{align*}

Combining the cases yields the explicit system:
\[
\widehat{\psi_a}(\mathbf{\beta})=
\begin{cases}
2^n-2s_a(2^m-1), & \text{if } \mathbf{\beta}=\mathbf{0},\\
2s_a, & \text{if } \mathbf{\beta}\notin V_\ell^\perp,\ \forall \; \ell\in \mathcal{S}_a,\\
-2^{m+1}+2s_a, & \text{if } \mathbf{\beta}\in V_\ell^\perp\setminus\{\mathbf{0}\} \text{ for some } \ell\in \mathcal{S}_a.
\end{cases}
\]
Substituting $\widehat{\psi_a}(\mathbf{\beta})$ directly into the weight equation $\mathrm{wt}(m_{(\psi_a)}(\mathbf{\beta})) = \frac{2^n - \widehat{\psi_a}(\mathbf{\beta})}{2}$ maps out the unique multiplicities in Table 1.
\end{proof}

\begin{table}[ht]
\centering
\caption{The values of Walsh transforms of $\phi \in \mathcal{F}$}
\label{Table_1}
\label{tab:walsh}
\renewcommand{\arraystretch}{1.0}
\begin{tabular}{|l|l|l|l|}
\hline
Walsh transform &
$\mathbf{\beta=0}$ &
$\mathbf{\beta} \notin V_\ell^\perp,\ \forall\, \ell \in \mathcal{S}$ &
$\mathbf{\beta}\in V_\ell^\perp\setminus\{\mathbf{0}\}\; \text{for some }\ell \in \mathcal{S}$ \\
\hline
$\widehat{\psi_a}(\mathbf{\beta})$ &
$2^n-2s_a(2^m-1)$ &
$2s_a$ &
$-2^{m+1}+2s_a$
\\
\hline
$\widehat{\psi_a+\psi_b}(\mathbf{\beta})$ &
$2^n-2\sigma_{ab}(2^m-1)$ &
$2\sigma_{ab}$ &
$-2^{m+1}+2\sigma_{ab}$
\\
\hline
$\widehat{\psi_a+\psi_b+\psi_c}(\mathbf{\beta})$ &
$2^n-2\sigma_{abc}(2^m-1)$ &
$2\sigma_{abc}$ &
$-2^{m+1}+2\sigma_{abc}$
\\
\hline
$\widehat{\displaystyle\sum_{k=1}^{4}\psi_k}(\mathbf{\beta})$ &
$2^n-2\sigma_{1234}(2^m-1)$ &
$2\sigma_{1234}$ &
$-2^{m+1}+2\sigma_{1234}$
\\
\hline
\end{tabular}
\end{table}

To guarantee that the linear code $\mathcal{M}_{(\psi_a)}$ is structurally minimal, we establish the following geometric and combinatorial constraints on the underlying index sets:

\begin{itemize}
    \item[\textbf{C1:}] $\mathcal{S}_c \not\subset \mathcal{S}_a \triangle \mathcal{S}_b$, $\mathcal{S}_a \triangle \mathcal{S}_b \not\subset \mathcal{S}_c$, $\mathcal{S}_d \not\subset \mathcal{S}_a \triangle \mathcal{S}_b \triangle \mathcal{S}_c$, and $\mathcal{S}_a \triangle \mathcal{S}_b \triangle \mathcal{S}_c \not\subset \mathcal{S}_d$ for  $a,b,c,d \in \{1,2,3,4\}$.
    \item[\textbf{C2:}] $\bigcap_{a=1}^3 \mathcal{S}_a \neq \emptyset$, $\bigcap_{a=1}^{4} \mathcal{S}_a \neq \emptyset$, and the condition $|\mathcal{S}_a \cap \mathcal{S}_b \cap \mathcal{S}_c| \neq |\mathcal{S}_i \cap \mathcal{S}_j|$ holds for at least two pairs of distinct indices and the condition $|\mathcal{S}_a \cap \mathcal{S}_b \cap \mathcal{S}_c \cap \mathcal{S}_d| \neq |\mathcal{S}_i \cap \mathcal{S}_j \cap \mathcal{S}_k|$ holds for at least three pairs of distinct indices.
    \item[\textbf{C3:}] $\sigma_{ab} \geq 2$, $\sigma_{abc} \geq 2$, $\sigma_{ab} + \sigma_{ad} \neq \sigma_{bd}$, $\sigma_{ab} + \sigma_{acd} \neq \sigma_{bcd}$ and $\sigma_{abc} + \sigma_{abd} \neq \sigma_{cd}$.
\end{itemize}

These constraints generalize the structural criteria developed for lower-dimensional configurations in \cite{LiuLiao2022}.

\begin{table}[h]
\centering
\caption{Weight distribution of $\mathcal{M}_{(\psi_a)}$}
\label{Table_2}
\begin{tabular}{|c|c|}
\hline
Weight ($\omega$) & Multiplicity ($A_\omega$) \\
\hline
$0$ & $1$ \\
\hline
$u(2^m-1)$ & $1, \quad u \in \{s_a, \sigma_{ab}, \sigma_{abc}, \sigma_{1234}\}$ \\
\hline
$2^{n-1}$ & $2^n-1$ \\
\hline
$2^{n-1}-u$ & $(2^m+1-u)(2^m-1), \quad u \in \{s_a, \sigma_{ab}, \sigma_{abc}, \sigma_{1234}\}$ \\
\hline
$2^{n-1}+2^m-u$ & $u(2^m-1), \quad u \in \{s_a, \sigma_{ab}, \sigma_{abc}, \sigma_{1234}\}$ \\
\hline
\end{tabular}
\end{table}

We present two fundamental properties derived from these set conditions, which are required to establish the minimality proofs of the code space.
\begin{lemma}\label{lem:set_theory}
For $a \in \{1,2,3,4\}$, let $\mathcal{S}_a$ be non-empty index sets. If constraint C1 holds, then:
\begin{enumerate}
    \item $\mathcal{S}_a \not\subseteq \mathcal{S}_b$ and $\mathcal{S}_a \cap \mathcal{S}_b \neq \emptyset$ for all distinct $a, b \in \{1,2,3,4\}$.
    \item $(\mathcal{S}_a \cap \mathcal{S}_b) \triangle (\mathcal{S}_a \cap \mathcal{S}_c) \subsetneq \mathcal{S}_b \triangle \mathcal{S}_c$ for distinct $a,b,c$.
\end{enumerate}
\end{lemma}

\begin{proof}
\begin{enumerate}
    \item Let $\mathcal{S}_a \subseteq \mathcal{S}_b$ then  $\mathcal{S}_a \triangle \mathcal{S}_b = \mathcal{S}_b \setminus \mathcal{S}_a$. so that 
    $\mathcal{S}_b \setminus \mathcal{S}_a \subseteq \mathcal{S}_b$ implies that $\mathcal{S}_a \triangle \mathcal{S}_b \subseteq \mathcal{S}_b$. This contradicts  C1. 
    Let $\mathcal{S}_a \cap \mathcal{S}_b = \emptyset$.
    
    Now if  $\mathcal{S}_a \triangle \mathcal{S}_b = \mathcal{S}_a \cup \mathcal{S}_b$  we have  $\mathcal{S}_a \subseteq \mathcal{S}_a \triangle \mathcal{S}_b$. Again, contradicts  C1.
    Thus, $\mathcal{S}_a \not\subseteq \mathcal{S}_b$ and $\mathcal{S}_a \cap \mathcal{S}_b \neq \emptyset$.

    \item Let $x \in (\mathcal{S}_a \cap \mathcal{S}_b) \triangle (\mathcal{S}_a \cap \mathcal{S}_c)$. Then 
    \begin{align*}
    &x \in (\mathcal{S}_a \cap \mathcal{S}_b) \setminus (\mathcal{S}_a \cap \mathcal{S}_c)\\ 
    &\implies x \in \mathcal{S}_a, \, x \in \mathcal{S}_b, \, x \notin \mathcal{S}_c\\
    & \implies x \in \mathcal{S}_b \setminus \mathcal{S}_c \subseteq \mathcal{S}_b \triangle \mathcal{S}_c\\
    & \implies (\mathcal{S}_a \cap \mathcal{S}_b) \triangle (\mathcal{S}_a \cap \mathcal{S}_c) \subseteq \mathcal{S}_b \triangle \mathcal{S}_c.
    \end{align*}
    Assume $(\mathcal{S}_a \cap \mathcal{S}_b) \triangle (\mathcal{S}_a \cap \mathcal{S}_c) = \mathcal{S}_b \triangle \mathcal{S}_c$.
    Then $(\mathcal{S}_a \cap \mathcal{S}_b) \triangle (\mathcal{S}_a \cap \mathcal{S}_c) \subseteq \mathcal{S}_a$ implies that  $\mathcal{S}_b \triangle \mathcal{S}_c \subseteq \mathcal{S}_a$,  this contradicts  C1. Thus, $(\mathcal{S}_a \cap \mathcal{S}_b) \triangle (\mathcal{S}_a \cap \mathcal{S}_c) \subsetneq \mathcal{S}_b \triangle \mathcal{S}_c$.
\end{enumerate}
\end{proof}

\begin{lemma}\label{lem:bound_theory}
Let $\mathcal{S}_a$ be non-empty index sets satisfying $2 \leq s_a \leq 2^{m-1}$ for $a \in \{1,2,3,4\}$. If conditions C1 and C2 hold simultaneously, the symmetric difference targets are upper-bounded by:
\[
\sigma_{ab} \leq 2^m - 2, \quad \quad \sigma_{abc} \leq 2^m - 2 \quad \text{and} \quad \sigma_{abcd} \leq 2^m - 2.
\]
\end{lemma}

\begin{proof}
By  Lemma \ref{lem:set_theory}(1) we have that  $\sigma_{ab} = s_a + s_b - 2s_{ab}, \quad s_{ab} \geq 1$. 
Given $s_a, s_b \leq 2^{m-1}$ it yields that $\sigma_{ab} \leq 2^{m-1} + 2^{m-1} - 2(1) = 2^m - 2$. 

Now $|\mathcal{S}_1 \cup \mathcal{S}_2 \cup \mathcal{S}_3| \leq 2^m + 1$. 
Expanding it gives  $\sigma_{abc} = |\mathcal{S}_1 \cup \mathcal{S}_2 \cup \mathcal{S}_3| - \left( s_{12} + s_{13} + s_{23} - 2s_{123} \right)$.
By Condition C2, $s_{12} + s_{13} + s_{23} - 2s_{123} \geq 3$. Therefore $\sigma_{abc} \leq (2^m + 1) - 3 = 2^m - 2$. 

Now, $|\mathcal{S}_1 \cup \mathcal{S}_2 \cup \mathcal{S}_3 \cup \mathcal{S}_4| \leq 2^m + 1$. Expanding it gives:
\[
\sigma_{abcd} = |\mathcal{S}_1 \cup \mathcal{S}_2 \cup \mathcal{S}_3 \cup \mathcal{S}_4| - \left( \sum_{1 \le a < b \le 4} s_{ab} - 3 \sum_{1 \le a < b < c \le 4} s_{abc} + 7 s_{abcd} \right).
\]
By Condition C2, $\sum_{1 \le a < b \le 4} s_{ab} - 3 \sum_{1 \le a < b < c \le 4} s_{abc} + 7 s_{abcd} \ge 3$. Therefore, $\sigma_{abcd} \le (2^m + 1) - 3 = 2^m - 2$.
\end{proof}



By utilizing the bounded discrete transform values from Lemma \ref{lem:bound_theory}, we can consolidate the multi-case propositions into two overarching theorems that satisfy the minimality criteria of Theorem \ref{Thm_Minimal_Cond}.




\begin{theorem}\label{Thm_Verify_Cond1}
Let $n \geq 8$ and $2 \leq s_a \leq 2^{m-1}$ for all $1 \leq a \leq 4$. For any distinct vectors $\mathbf{x}, \mathbf{y} \in \mathbb{F}_2^n$, the Walsh-Hadamard combinations of the spread functions strictly satisfy:
\[ \widehat{\phi}_1(\mathbf{x}) + \widehat{\phi}_2(\mathbf{y}) \neq 2^n \quad \text{and} \quad \widehat{\phi}_1(\mathbf{x}) - \widehat{\phi}_2(\mathbf{y}) \neq 2^n, \]
for all choices of $\phi_1, \phi_2 \in \mathcal{F}$.
\end{theorem}

\begin{proof}
First, we have $\max_{\mathbf{\beta} \neq \mathbf{0}} |\widehat{\phi}(\mathbf{\beta})| \leq 2^{m+1} - 4$. Let $\sigma_1 = |\mathcal{S}_{\phi_1}|$ and $\sigma_2 = |\mathcal{S}_{\phi_2}|$. By Lemma \ref{lem:bound_theory}:
\begin{equation}\label{eq:thm_bounds}
2 \leq \sigma_1, \sigma_2 \leq 2^m - 2.
\end{equation}
We have the following cases:

\textbf{Case 1: $\mathbf{x} \neq \mathbf{0}$ and $\mathbf{y} \neq \mathbf{0}$.} \\
By triangle inequality, $|\widehat{\phi}_1(\mathbf{x}) \pm \widehat{\phi}_2(\mathbf{y})| \leq |\widehat{\phi}_1(\mathbf{x})| + |\widehat{\phi}_2(\mathbf{y})| \leq 2(2^{m+1}-4) = 2^{m+2}-8$, which  implies $\widehat{\phi}_1(\mathbf{x}) \pm \widehat{\phi}_2(\mathbf{y}) \neq 2^n$.


\textbf{Case 2: $\mathbf{x} = \mathbf{0}$ and $\mathbf{y} \neq \mathbf{0}$.} \\
Substituting $\mathbf{x} = \mathbf{0}$ into $\psi_a(\mathbf{x})=\sum_{\mathbf{x} \in \mathbb{F}_2^n} (-1)^{\mathbf{\beta} \cdot \mathbf{x}} - 2 \sum_{\ell \in \mathcal{S}_a} \sum_{\mathbf{x} \in V_\ell} (-1)^{\mathbf{\beta} \cdot \mathbf{x}} + 2s_a$  yields
that 
$\widehat{\phi}_1(\mathbf{0}) \pm \widehat{\phi}_2(\mathbf{y}) = 2^n - 2\sigma_1(2^m - 1) \pm \widehat{\phi}_2(\mathbf{y})$.
Now assume $\widehat{\phi}_1(\mathbf{0}) + \widehat{\phi}_2(\mathbf{y}) = 2^n$. This implies
\begin{equation}\label{eq:case2_pos_eq}
\widehat{\phi}_2(\mathbf{y}) = 2\sigma_1(2^m - 1).
\end{equation}
Since $\mathbf{y} \in V_{k_0}^\perp \setminus \{\mathbf{0}\}$ for at most one $k_0 \in \mathcal{S}_{\phi_2}$, we have 
\begin{equation}
\max_{\mathbf{y} \neq \mathbf{0}} \widehat{\phi}_2(\mathbf{y}) \leq 2\sigma_2 + 2^{m+1}. 
\end{equation}
Now, this together  with \eqref{eq:case2_pos_eq} gives 
$2\sigma_2 + 2^{m+1} \geq 2\sigma_1(2^m - 1)$ if and only if $\sigma_2 + 2^m \geq \sigma_1 2^m - \sigma_1$.

Applying \eqref{eq:thm_bounds} we have $(2^m - 2) + 2^m \geq 2^{m+1} - 2$ if and only if $(\sigma_1, \sigma_2) = (2, 2^m - 2)$.  By conditions C1 and C3, $\mathcal{S}_{\phi_1} \not\subseteq \mathcal{S}_{\phi_2}$ and $\mathcal{S}_{\phi_2} \not\subseteq \mathcal{S}_{\phi_1}$ implies $\sigma_2 \leq 2^m - 3$ when $\sigma_1 = 2$, which results in a contradiction.

Now assume $\widehat{\phi}_1(\mathbf{0}) - \widehat{\phi}_2(\mathbf{y}) = 2^n$. This implies that 
$-\widehat{\phi}_2(\mathbf{y}) = 2\sigma_1(2^m - 1)$.
Since $\sigma_1 \geq 2$, we have $2\sigma_1(2^m - 1) \geq 4(2^m - 1) = 2^{m+2} - 4$. However:
\[
\max_{\mathbf{y} \neq \mathbf{0}} |-\widehat{\phi}_2(\mathbf{y})| \leq 2^{m+1} - 2\sigma_2 < 2^{m+2} - 4,
\]
which is impossible.

By symmetry, the case $\mathbf{x} \neq \mathbf{0}, \mathbf{y} = \mathbf{0}$ holds identically. Thus, $\widehat{\phi}_1(\mathbf{x}) \pm \widehat{\phi}_2(\mathbf{y}) \neq 2^n$ for all $\mathbf{x} \neq \mathbf{y}$.
\end{proof}




\begin{theorem}\label{Thm_Verify_Cond2}
Let $n \geq 8$ and $2 \leq s_a \leq 2^{m-1}$ for all $1 \leq a \leq 4$. For all functional pairings $\phi_1, \phi_2 \in \mathcal{F}$ with $\phi_1 \neq \phi_2$, the shifting differences satisfy:
\[ \widehat{\phi_1}(\mathbf{x}+\mathbf{y})+\widehat{\phi_2}(\mathbf{x})-\widehat{\phi_1+\phi_2}(\mathbf{y})\neq 2^n, \]
across all vector spatial configurations of $\mathbf{x}, \mathbf{y} \in \mathbb{F}_2^n$.
\end{theorem}

\begin{proof}
Let $\phi_1, \phi_2 \in \mathcal{F}$ with $\phi_1 \neq \phi_2$, and define $\psi = \phi_1 + \phi_2 \in \mathcal{F}$. Let $\sigma_1 = |\mathcal{S}_{\phi_1}|$, $\sigma_2 = |\mathcal{S}_{\phi_2}|$, and $\sigma_3 = |\mathcal{S}_{\psi}| = |\mathcal{S}_{\phi_1} \triangle \mathcal{S}_{\phi_2}|$. For any $\mathbf{x}, \mathbf{y} \in \mathbb{F}_2^n$, define $\Lambda(\mathbf{x}, \mathbf{y}) = \widehat{\phi}_1(\mathbf{x}+\mathbf{y}) + \widehat{\phi}_2(\mathbf{x}) - \widehat{\psi}(\mathbf{y})$.


We evaluate $\Lambda(\mathbf{x}, \mathbf{y})$ across three disjoint spatial vector configurations:

\textbf{Case 1: $\mathbf{x} = \mathbf{y}$.} \\
When $\mathbf{x} = \mathbf{y}$,  $\Lambda(\mathbf{x}, \mathbf{x}) = \widehat{\phi}_1(\mathbf{0}) + \widehat{\phi}_2(\mathbf{x}) - \widehat{\psi}(\mathbf{x})$. Substituting the evaluation at the origin yields:
\[
\Lambda(\mathbf{x}, \mathbf{x}) = 2^n - 2\sigma_1(2^m - 1) + \widehat{\phi}_2(\mathbf{x}) - \widehat{\psi}(\mathbf{x}).
\]
Assuming $\Lambda(\mathbf{x}, \mathbf{x}) = 2^n$ requires $\widehat{\phi}_2(\mathbf{x}) - \widehat{\psi}(\mathbf{x}) = 2\sigma_1(2^m - 1)$. If $\mathbf{x} = \mathbf{0}$, this gives:
\[
2^n - 2\sigma_2(2^m - 1) - \left(2^n - 2\sigma_3(2^m - 1)\right) = 2\sigma_1(2^m - 1) \iff \sigma_3 = \sigma_1 + \sigma_2,
\]
which contradicts $|\mathcal{S}_{\phi_1} \triangle \mathcal{S}_{\phi_2}| \leq \sigma_1 + \sigma_2 - 2|\mathcal{S}_{\phi_1} \cap \mathcal{S}_{\phi_2}| < \sigma_1 + \sigma_2$ by condition C2 ($|\mathcal{S}_{\phi_1} \cap \mathcal{S}_{\phi_2}| \neq 0$). If $\mathbf{x} \neq \mathbf{0}$, then $\max_{\mathbf{x} \neq \mathbf{0}} |\widehat{\phi}_2(\mathbf{x}) - \widehat{\psi}(\mathbf{x})| \leq 2^{m+2} - 8$, whereas $2\sigma_1(2^m - 1) \geq 2^{m+2} - 4$, making equality impossible.

\textbf{Case 2: Either $\mathbf{x} = \mathbf{0}$ or $\mathbf{y} = \mathbf{0}$ with $\mathbf{x} \neq \mathbf{y}$.} \\
If $\mathbf{x} = \mathbf{0}$ and $\mathbf{y} \neq \mathbf{0}$, then $\Lambda(\mathbf{0}, \mathbf{y}) = \widehat{\phi}_1(\mathbf{y}) + \widehat{\phi}_2(\mathbf{0}) - \widehat{\psi}(\mathbf{y})$. Assume $\Lambda(\mathbf{0}, \mathbf{y}) = 2^n$:
\[
\widehat{\phi}_1(\mathbf{y}) - \widehat{\psi}(\mathbf{y}) = 2\sigma_2(2^m - 1).
\]
By the spectral triangle inequality:
\[
\max_{\mathbf{y} \neq \mathbf{0}} |\widehat{\phi}_1(\mathbf{y}) - \widehat{\psi}(\mathbf{y})| \leq (2\sigma_1 + 2^{m+1}) + (2\sigma_3 + 2^{m+1}) = 2(\sigma_1 + \sigma_3) + 2^{m+2}.
\]
This forces $2(\sigma_1 + \sigma_3) + 2^{m+2} \geq 2\sigma_2(2^m - 1)$, which yields a direct contradiction under constraints C1–C3. The symmetric case $\mathbf{y} = \mathbf{0}, \mathbf{x} \neq \mathbf{0}$ holds identically.

\textbf{Case 3: $\mathbf{x} \neq \mathbf{0}$, $\mathbf{y} \neq \mathbf{0}$, and $\mathbf{x} \neq \mathbf{y}$.} \\
Since $\mathbf{x}, \mathbf{y}, \mathbf{x}+\mathbf{y}$ are all non-zero, each vector belongs to at most one orthogonal subspace component $V_\ell^\perp$. By the triangle inequality:
\[
|\Lambda(\mathbf{x}, \mathbf{y})| \leq |\widehat{\phi}_1(\mathbf{x}+\mathbf{y})| + |\widehat{\phi}_2(\mathbf{x})| + |\widehat{\psi}(\mathbf{y})| \leq 3(2^{m+1} - 4) = 3 \cdot 2^{m+1} - 12.
\]
For $n \geq 8 \implies m \geq 4$:
\[
3 \cdot 2^{m+1} - 12 < 2^{2m} = 2^n.
\]
Thus, $\Lambda(\mathbf{x}, \mathbf{y}) \neq 2^n$ across all vector spatial configurations.
\end{proof}
\begin{theorem}\label{Thm_Main_Code_Result}
Let $n \ge 8$ be an even integer with $m = n/2$, and let the index set cardinalities satisfy $2 \le s_a \le 2^{m-1}$ for all $1 \le a \le 4$. Define the minimum partial spread intersection parameter across all 15 non-zero functions $\phi \in \mathcal{F}$ as:
\[
\theta = \min_{1 \le a,b,c \le 4} \{ s_a, \, \sigma_{ab}, \, \sigma_{abc}, \, \sigma_{1234} \}.
\]
If $\theta \le 2^{m-1} - 1$, then $\mathcal{M}_{(\psi_a)}$ is a valid $[2^n - 1, \, n + 4, \, \theta(2^m - 1)]$ minimal binary linear code. Furthermore, if $\theta \le 2^{m-2}$, the code strictly violates the Ashikhmin-Barg condition, exhibiting a Hamming weight ratio of:
\[
\frac{w_{\min}}{w_{\max}} \le \frac{1}{2}.
\]
\end{theorem}
\begin{proof}
By Theorems \ref{Thm_Verify_Cond1} and \ref{Thm_Verify_Cond2}, all spatial vector conditions required by Theorem \ref{Thm_Minimal_Cond} are simultaneously satisfied, establishing that every non-zero codeword in $\mathcal{M}_{(\psi_a)}$ is minimal.

From the weight spectrum derived in Table~\ref{Table_2}, the non-zero Hamming weights of $\mathcal{M}_{(\psi_a)}$ belong to the set:
\[
\Omega \setminus \{0\} = \{ u(2^m - 1), \; 2^{n-1}, \; 2^{n-1} - u, \; 2^{n-1} + 2^m - u \},
\]
where $u \in \{s_a, \sigma_{ab}, \sigma_{abc}, \sigma_{1234}\}$. The absolute minimum weight of $\mathcal{M}_{(\psi_a)}$ is attained at $u = \theta$, giving $w_{\min} = \theta(2^m - 1)$.

Conversely, the maximum weight is attained at $w_{\max} = 2^{n-1} + 2^m - \theta$. Evaluating the ratio $w_{\min} / w_{\max}$ under the structural threshold $\theta \le 2^{m-2}$ gives:
\[
\frac{w_{\min}}{w_{\max}} = \frac{\theta(2^m - 1)}{2^{n-1} + 2^m - \theta}.
\]
Since $n \ge 8$ implies $m \ge 4$, substituting $\theta \le 2^{m-2}$ yields:
\[
\frac{w_{\min}}{w_{\max}} \le \frac{2^{m-2}(2^m - 1)}{2^{2m-1} + 2^m - 2^{m-2}} = \frac{2^{2m-2} - 2^{m-2}}{2^{2m-1} + 3 \cdot 2^{m-2}} < \frac{1}{2}.
\]
Thus, the code strictly violates the Ashikhmin--Barg ceiling ($w_{\min} / w_{\max} \le 1/2$) while remaining strictly minimal.
\end{proof}

\begin{example}\label{ex:AB_strictly_greater}
Let $n = 8$ ($m = 4$), yielding code length $N = 2^8 - 1 = 255$ and dimension $K = n + 4 = 12$. We select $17$ pairwise disjoint $4$-dimensional subspaces $\{V_1, V_2, \ldots, V_{17}\}$ in $\mathbb{F}_2^8$. To ensure conditions C1--C3 and the size bounds $2 \le s_a \le 2^{m-1} = 8$ are strictly satisfied, we define the four basis index sets as:
\[
\mathcal{S}_1 = \{1, 2, 3, 4, 10, 11, 13\}, \quad \mathcal{S}_2 = \{3, 4, 5, 6, 10, 11, 13\},
\]
\[
\mathcal{S}_3 = \{1, 3, 7, 8, 10, 12, 13\}, \quad \text{and} \quad \mathcal{S}_4 = \{2, 3, 5, 9, 11, 12, 13\},
\]
giving individual set cardinalities $s_1 = s_2 = s_3 = s_4 = 7 \le 2^{4-1}$.
Evaluating all $15$ non-zero function combinations in $\mathcal{F}$ yields the parameter set:
\[
u \in \{\sigma_{ab}, \sigma_{abc}, \sigma_{1234}\} \subset \{5, 6, 7, 8\},
\]
thereby determining the absolute minimum partial spread intersection parameter $\theta = \min \{s_a, \sigma_{ab}, \sigma_{abc}, \sigma_{1234}\} = 5$.

By Theorem \ref{Thm_Main_Code_Result}, the resulting code $\mathcal{M}_{(\psi_a)}$ is a minimal binary linear code with parameters $[255, 12, 75]$, whose exact weight enumerator polynomial is given by:
\[
\begin{aligned}
W(z) = 1 &+ z^{75} + 5z^{90} + 5z^{105} + 544z^{120} + 675z^{121} + 825z^{122} + 180z^{123} \\
&+ 255z^{128} + 480z^{136} + 525z^{137} + 450z^{138} + 75z^{139}.
\end{aligned}
\] From Table 2, the absolute minimum and maximum Hamming weights evaluate to $w_{\min} = 75$ and $w_{\max} = 140$, respectively. Evaluating their weight ratio gives:
\[
\frac{w_{\min}}{w_{\max}} = \frac{75}{140} > \frac{1}{2}.
\]
Thus, because $\theta = 5 > 2^{m-2} = 4$, this parameter instance directly yields a minimal linear code that \textbf{strictly satisfies the classical Ashikhmin--Barg sufficient condition} ($\frac{w_{\min}}{w_{\max}} > \frac{1}{2}$), serving as an immediate baseline application of Theorem \ref{Thm_Main_Code_Result}. 

\end{example}


\begin{example}\label{ex:AB_violated}
Let $n = 8$ ($m = 4$), yielding code length $N = 2^8 - 1 = 255$ and dimension $K = n + 4 = 12$. We select $17$ pairwise disjoint $4$-dimensional subspaces $\{V_1, V_2, \ldots, V_{17}\}$ in $\mathbb{F}_2^8$. To ensure conditions C1--C3 and the size bounds $2 \le s_a \le 2^{m-1} = 8$ are strictly satisfied, we define the four basis index sets as:
\[
\mathcal{S}_1 = \{1, 2, 3, 4, 10, 11\}, \quad \mathcal{S}_2 = \{3, 4, 5, 6, 10, 11\},
\]
\[
\mathcal{S}_3 = \{1, 3, 7, 8, 10, 12\}, \quad \text{and} \quad \mathcal{S}_4 = \{2, 3, 5, 9, 11, 12\},
\]
giving individual set cardinalities $s_1 = s_2 = s_3 = s_4 = 6 \le 2^{4-1}$.

Evaluating all $15$ non-zero function combinations in $\mathcal{F}$ yields the parameter set:
\[
u \in \{\sigma_{ab}, \sigma_{abc}, \sigma_{1234}\} \subset \{4, 5, 6, 7, 8\},
\]
thereby determining the absolute minimum partial spread intersection parameter $\theta = \min \{s_a, \sigma_{ab}, \sigma_{abc}, \sigma_{1234}\} = 4$ (attained at $\sigma_{12} = 4$).

By Theorem \ref{Thm_Main_Code_Result}, the resulting code $\mathcal{M}_{(\psi_a)}$ is a minimal binary linear code with parameters $[255, 12, 60]$, whose exact weight enumerator polynomial is given by:
\[
\begin{aligned}
W(z) = 1 &+ z^{60} + 5z^{90} + 5z^{105} + 544z^{120} + 750z^{121} + 825z^{122} + 195z^{124} \\
&+ 255z^{128} + 480z^{136} + 525z^{137} + 450z^{138} + 60z^{140}.
\end{aligned}
\]
Here, the absolute minimum and maximum Hamming weights evaluate to $w_{\min} = 60$ and $w_{\max} = 140$, respectively. Evaluating their weight ratio gives:
\[
\frac{w_{\min}}{w_{\max}} = \frac{60}{140}< \frac{1}{2}.
\]
Thus, because $\theta = 4 \le 2^{m-2} = 4$, this parameter instance directly demonstrates a minimal linear code that \textbf{strictly violates the classical Ashikhmin--Barg sufficient condition} ($\frac{w_{\min}}{w_{\max}} \le \frac{1}{2}$), highlighting the geometric sieving capability of Theorem \ref{Thm_Main_Code_Result}.
\end{example}







\section{Applications in Cryptographic Secret Sharing Schemes}

Let $\mathcal{M}$ be an $[N, K, D]$ binary linear code with generator matrix $\mathbf{G} = [\mathbf{g}_0, \mathbf{g}_1, \ldots, \mathbf{g}_{N-1}] \in \mathbb{F}_2^{K \times N}$. In a Massey secret sharing scheme, a trusted dealer distributes a secret $s_0 = \mathbf{u} \cdot \mathbf{g}_0 \in \mathbb{F}_2$ across $N-1$ participants $\mathcal{P} = \{P_1, P_2, \ldots, P_{N-1}\}$ by assigning shares $s_i = \mathbf{u} \cdot \mathbf{g}_i$ for $1 \le i \le N-1$, where $\mathbf{u} \in \mathbb{F}_2^K$ is chosen uniformly at random.

A participant coalition $\mathcal{P}_A = \{P_{i_1}, P_{i_2}, \ldots, P_{i_m}\} \subseteq \mathcal{P}$ is called a \emph{minimal access set} if and only if $\text{rank}([\mathbf{g}_0, \mathbf{g}_{i_1}, \ldots, \mathbf{g}_{i_m}]) = \text{rank}([\mathbf{g}_{i_1}, \ldots, \mathbf{g}_{i_m}])$, with no proper subset satisfying this rank condition. Let $\Gamma$ denote the collection of all minimal access sets. By Massey \cite{Massey1993}, $\mathcal{P}_A \in \Gamma$ if and only if there exists a minimal codeword $\mathbf{c} = (c_0, c_1, \ldots, c_{N-1}) \in \mathcal{M}^\perp$ such that $c_0 = 1$ and $\mathrm{Supp}(\mathbf{c}) = \{0\} \cup \{i_j \mid P_{i_j} \in \mathcal{P}_A\}$.

When $\mathcal{M}$ is a minimal linear code, every non-zero codeword is minimal, guaranteeing that the dual code $\mathcal{M}^\perp$ induces an \emph{ideal} secret sharing scheme where $\Gamma$ is completely characterized by the affine cross-section $c_0 = 1$.



\subsection{Cryptographic Significance of the $n+4$ Access Architecture}
 In this section, we apply the newly constructed family of minimal binary linear codes $\mathcal{M}_{(\psi_a)}$ of dimension $K = n+4$ to design an ideal secret sharing scheme using Massey's foundational framework \cite{Massey1993}.

To formally establish the structural foundation of the access structure induced by our $K = n+4$ code family, we first determine the exact cardinality and support identity of the authorized minimal access sets.

\begin{theorem}\label{Thm:SSS_Construction}
Let $\mathcal{M}_{(\psi_a)}$ be the $[2^n-1, n+4, \theta(2^m-1)]$ minimal binary linear code constructed in Theorem \ref{Thm_Main_Code_Result} over length $N = 2^n - 1$ and dimension $K = n+4$. The access structure $\Gamma$ over $\mathcal{P} = \{P_1, \ldots, P_{2^n-2}\}$ satisfies:
\begin{enumerate}
    \item  $|\Gamma| = 2^{n+3}$.
    \item  For every $\mathcal{P}_A \in \Gamma$, the cardinality $|\mathcal{P}_A|$ is uniquely given by:
    \[
    |\mathcal{P}_A| = w - 1, \quad \text{for some } w \in \Omega \setminus \{0\},
    \]
    where $\Omega$ is the weight spectrum of $\mathcal{M}_{(\psi_a)}$ established in Table 2.
\end{enumerate}
\end{theorem}

\begin{proof}
(i) Define the affine linear subspace $\mathcal{M}_1 = \{\mathbf{c} \in \mathcal{M}_{(\psi_a)} \mid c_0 = 1\}$. Since $\mathcal{M}_{(\psi_a)}$ contains non-zero codewords with non-zero first coordinate, $|\mathcal{M}_1| = 2^{K-1} = 2^{(n+4)-1} = 2^{n+3}$. By Theorem \ref{Thm_Main_Code_Result}, every $\mathbf{c} \in \mathcal{M}_{(\psi_a)} \setminus \{\mathbf{0}\}$ is minimal. Thus, every $\mathbf{c} \in \mathcal{M}_1$ uniquely induces a minimal access set $\mathcal{P}_A = \{P_i \mid c_i = 1, i \neq 0\} \in \Gamma$, establishing $|\Gamma| = 2^{n+3}$.

(ii) For any $\mathbf{c} \in \mathcal{M}_1$, its Hamming weight is $wt(\mathbf{c}) = |\mathrm{Supp}(\mathbf{c})| = 1 + |\mathcal{P}_A|$. Thus, $|\mathcal{P}_A| = wt(\mathbf{c}) - 1$. The values $wt(\mathbf{c})$ are restricted to the non-zero weight support $\Omega \setminus \{0\}$ of $\mathcal{M}_{(\psi_a)}$.
\end{proof}

Having established the exact total number of authorized paths, we now quantify the operational boundaries and overall size distribution of these access sets.

\begin{proposition}\label{prop:access_bounds}
Let $\theta = \min\{s_a, \sigma_{ab}, \sigma_{abc}, \sigma_{1234}\}$ and $u_{\max} = \max\{s_a, \sigma_{ab}, \sigma_{abc}, \sigma_{1234}\} \le 2^m - 2$. The cardinality of any minimal access set $\mathcal{P}_A \in \Gamma$ satisfies:
\[
2^{n-1} - 1 - u_{\max} \le |\mathcal{P}_A| \le 2^{n-1} - 1 + 2^m - \theta.
\]
Moreover, if $\theta \le 2^{m-2}$, the operational span $\Delta = \max_{\mathcal{P}_A}|\mathcal{P}_A| - \min_{\mathcal{P}_A}|\mathcal{P}_A|$ satisfies the lower bound:
\[
\Delta \ge 2^{n-1} + 2^{m-2} - 1,
\]
establishing a multi-threshold access spectrum.
\end{proposition}

\begin{proof}
From Table 2, the non-zero weight values of $\mathcal{M}_{(\psi_a)}$ belong to the set:
\[
\Omega \setminus \{0\} = \{u(2^m-1), \, 2^{n-1}, \, 2^{n-1}-u, \, 2^{n-1}+2^m-u\},
\]
where $u \in \{s_a, \sigma_{ab}, \sigma_{abc}, \sigma_{1234}\}$. The absolute minimum weight in $\mathcal{M}_{(\psi_a)}$ is $w_{\min} = \theta(2^m-1)$, and the maximum weight is $w_{\max} = 2^{n-1} + 2^m - \theta$. Applying $|\mathcal{P}_A| = w - 1$ gives the extrema:
\[
\min_{\mathcal{P}_A \in \Gamma} |\mathcal{P}_A| = 2^{n-1} - 1 - u_{\max}, \quad \text{and} \quad \max_{\mathcal{P}_A \in \Gamma} |\mathcal{P}_A| = 2^{n-1} - 1 + 2^m - \theta.
\]
The operational span evaluates to:
\[
\Delta = (2^{n-1} - 1 + 2^m - \theta) - (2^{n-1} - 1 - u_{\max}) = 2^m - \theta + u_{\max}.
\]
Substituting $\theta \le 2^{m-2}$ and $u_{\max} \ge 2^{m-1}$ into $\Delta$ yields $\Delta \ge 2^m - 2^{m-2} + 2^{m-1} = 2^{n-1} + 2^{m-2} - 1$.
\end{proof}

\begin{table}[htbp]
\centering
\small
\caption{Comparison of Secret Sharing Schemes Derived from Partial Spread Minimal Codes}
\label{tab:SSS_Comparison}
\begin{tabular}{|c|c|c|c|}
\hline
  \textbf{Dimension ($K$)} & \textbf{Access Sets ($|\Gamma|$)} & \textbf{AB Condition Violation} & \textbf{Operational Span ($\Delta$)} \\ \hline
  $n+1, n+2$ & $2^{n+1}$\cite{LiuLiao2022} & Limited / Restrictive & Narrow ($\approx 2^{m-1}$) \\ \hline
 $n+2$ & $2^{n+1}$\cite{LiYue2020} & Satisfied & Moderate \\ \hline
  $n+2, n+3$ & $2^{n+2}$\cite{MesnagerQianCaoYuan2023} & Partial & Medium \\ \hline
$\mathbf{n+4}$ & $\mathbf{2^{n+3}}$ & \textbf{Strictly Violated ($\le \frac{1}{2}$)} & \textbf{Broad ($\ge 2^{n-1}+2^{m-2}-1$)} \\ \hline
\end{tabular}
\end{table}



As highlighted in Table \ref{tab:SSS_Comparison}, the secret sharing schemes induced by our $[2^n-1, n+4]$ minimal code family exhibit several fundamental cryptographic advantages over existing constructions in the literature \cite{LiYue2020, LiuLiao2022, MesnagerQianCaoYuan2023}:
\begin{itemize}
    \item \textbf{Expanded Authorization Density:} Over the same block length $N = 2^n-1$, our $K = n+4$ construction expands the total number of minimal access sets to $|\Gamma| = 2^{n+3}$, quadrupling the authorization pathways available in $(n+2)$-dimensional schemes \cite{LiuLiao2022} and doubling those in $(n+3)$-dimensional schemes \cite{MesnagerQianCaoYuan2023}.
    \item \textbf{Multi-Threshold Flexibility via Structural AB Violation:} While classical schemes constrained by the Ashikhmin--Barg condition \cite{AshikhminBarg1998, CarletDingYuan2005} restrict access set cardinalities to narrow clusters, our construction strictly achieves $w_{\min}/w_{\max} \le 1/2$. This yields a multi-threshold access spectrum with an operational span $\Delta \ge 2^{n-1} + 2^{m-2} - 1$, providing ideal multi-tiered authorization capabilities for hierarchical multi-party computation.
    \item \textbf{Throughput Efficiency:} By Proposition 4.3, our scheme achieves an asymptotic relative information rate enhancement of $\frac{2}{n+2}$ over standard $(n+2)$-dimensional constructions, translating to a $20\%$ throughput gain for $n=8$.
\end{itemize}



To highlight the efficiency gains of expanding the dimension to $K = n+4$ relative to lower-dimensional partial spread codes, we analyze the asymptotic code rate and relative information gain.

\begin{proposition}\label{prop:information_rate}
Let $\mathcal{M}_{(\psi_a)}$ be the $[2^n-1, n+4, \theta(2^m-1)]$ minimal binary linear code constructed in Theorem \ref{Thm_Main_Code_Result} over length $N = 2^n - 1$. The information rate $R$ of the code and the access density ratio $\rho = \frac{\log_2 |\Gamma|}{N}$ satisfy the asymptotic limits:
\[
\lim_{n \to \infty} R = \lim_{n \to \infty} \frac{n+4}{2^n - 1} = 0, \quad \text{and} \quad \lim_{n \to \infty} \frac{\rho}{R} = 1.
\]
Furthermore, for any dimension $n \ge 8$, the relative information rate enhancement over the standard $(n+2)$-dimensional construction \cite{LiuLiao2022} satisfies:
\[
\frac{R_{n+4} - R_{n+2}}{R_{n+2}} = \frac{2}{n+2},
\]
yielding a strictly positive throughput gain over dual secret sharing schemes.
\end{proposition}

\begin{proof}
By Theorem \ref{Thm_Main_Code_Result}, $K = n+4$ and $N = 2^n - 1$. The code rate is $R = \frac{n+4}{2^n - 1}$. By Theorem \ref{Thm:SSS_Construction}, the number of minimal access sets is $|\Gamma| = 2^{n+3}$. The access density is given by $\rho = \frac{\log_2(2^{n+3})}{2^n - 1} = \frac{n+3}{2^n - 1}$. Taking the ratio yields $\frac{\rho}{R} = \frac{n+3}{n+4} = 1 - \frac{1}{n+4}$. Taking $n \to \infty$, we obtain $\lim_{n \to \infty} \frac{\rho}{R} = 1$. The relative dimension improvement over an $(n+2)$-dimensional code evaluates to $\frac{(n+4)/(2^n-1) - (n+2)/(2^n-1)}{(n+2)/(2^n-1)} = \frac{2}{n+2}$.
\end{proof}

From a cryptographic security perspective, it is necessary to examine how the system resists unauthorized collusion and statistical share approximation. The following result provides an exact bound on cheating probabilities and structural variance.

\begin{theorem}\label{thm:cheating_prob}
In the Massey secret sharing scheme induced by $\mathcal{M}_{(\psi_a)}^\perp$, let $\mathcal{P}_C \subset \mathcal{P}$ be an unauthorized coalition of participants such that $\mathcal{P}_C \notin \Gamma$ and $|\mathcal{P}_C| = k_c < \min_{\mathcal{P}_A \in \Gamma} |\mathcal{P}_A|$. The probability $P_{\mathrm{cheat}}$ that $\mathcal{P}_C$ successfully reconstructs the secret $s_0$ without authorized shares satisfies:
\[
P_{\mathrm{cheat}} = \frac{1}{2}.
\]
Furthermore, the statistical variance $\mathrm{Var}(|\mathcal{P}_A|)$ of the access set cardinalities across $\Gamma$ strictly satisfies:
\[
\mathrm{Var}(|\mathcal{P}_A|) > 0,
\]
rendering the access structure immune to deterministic weight approximation attacks.
\end{theorem}

\begin{proof}
Since $\mathcal{M}_{(\psi_a)}$ is a minimal linear code, every non-zero codeword in $\mathcal{M}_{(\psi_a)}^\perp$ is a minimal support vector. For any unauthorized coalition $\mathcal{P}_C$ with $k_c < \min |\mathcal{P}_A|$, the submatrix $\mathbf{G}_{\mathcal{P}_C}$ formed by the columns of $\mathbf{G}$ corresponding to $\mathcal{P}_C$ satisfies $\mathbf{g}_0 \notin \text{span}(\mathbf{G}_{\mathcal{P}_C})$. Consequently, the secret $s_0 = \mathbf{u} \cdot \mathbf{g}_0$ is statistically independent of the coalition shares $\mathbf{s}_{\mathcal{P}_C} = \mathbf{u} \cdot \mathbf{G}_{\mathcal{P}_C}$. Over $\mathbb{F}_2$, the conditional probability evaluates to $P(s_0 = \alpha \mid \mathbf{s}_{\mathcal{P}_C}) = \frac{1}{2}$ for all $\alpha \in \mathbb{F}_2$, establishing $P_{\mathrm{cheat}} = 1/2$.

To evaluate the variance $\mathrm{Var}(|\mathcal{P}_A|)$, note from Table 2 that the non-zero weight set $\Omega \setminus \{0\}$ contains at least four distinct values when $\theta \le 2^{m-2}$. Since $|\mathcal{P}_A| = w - 1$ assumes distinct integer values with non-zero multiplicities $A_w > 0$, the cardinalities do not collapse to a single constant, which strictly forces $\mathrm{Var}(|\mathcal{P}_A|) = E[|\mathcal{P}_A|^2] - (E[|\mathcal{P}_A|])^2 > 0$.
\end{proof}

Next, we analyze the structural overlap between authorized groups to show that no two distinct coalitions can share a dominant majority of participants.

\begin{proposition}\label{prop:intersection_hierarchy}
Let $\mathcal{P}_{A_1}, \mathcal{P}_{A_2} \in \Gamma$ be two distinct minimal access sets corresponding to dual codewords $\mathbf{c}_1, \mathbf{c}_2 \in \mathcal{M}_1 = \{\mathbf{c} \in \mathcal{M}_{(\psi_a)} \mid c_0 = 1\}$. The cardinality of their participant overlap $\mathcal{P}_{A_1} \cap \mathcal{P}_{A_2}$ satisfies:
\[
|\mathcal{P}_{A_1} \cap \mathcal{P}_{A_2}| = \frac{1}{2} \left( |\mathcal{P}_{A_1}| + |\mathcal{P}_{A_2}| - wt(\mathbf{c}_1 + \mathbf{c}_2) \right).
\]
In particular, the maximum participant overlap is strictly upper-bounded by:
\[
|\mathcal{P}_{A_1} \cap \mathcal{P}_{A_2}| \le 2^{n-1} + 2^m - \theta - 1 - 2^{m-1} \cdot \theta.
\]
\end{proposition}

\begin{proof}
Let $\mathbf{c}_1, \mathbf{c}_2 \in \mathcal{M}_1$. By definition, $c_{1,0} = c_{2,0} = 1$. The support sets of the corresponding access sets are $\text{Supp}(\mathbf{c}_1) = \{0\} \cup \mathcal{P}_{A_1}$ and $\text{Supp}(\mathbf{c}_2) = \{0\} \cup \mathcal{P}_{A_2}$. The sum $\mathbf{c}_3 = \mathbf{c}_1 + \mathbf{c}_2 \in \mathcal{M}_{(\psi_a)}$ has a zero first coordinate ($c_{3,0} = 1 + 1 = 0$). By standard binary support expansion:
\[
wt(\mathbf{c}_1 + \mathbf{c}_2) = wt(\mathbf{c}_1) + wt(\mathbf{c}_2) - 2 |\mathrm{Supp}(\mathbf{c}_1) \cap \mathrm{Supp}(\mathbf{c}_2)|.
\]
Since $\mathrm{Supp}(\mathbf{c}_1) \cap \mathrm{Supp}(\mathbf{c}_2) = \{0\} \cup (\mathcal{P}_{A_1} \cap \mathcal{P}_{A_2})$, we have $|\mathrm{Supp}(\mathbf{c}_1) \cap \mathrm{Supp}(\mathbf{c}_2)| = 1 + |\mathcal{P}_{A_1} \cap \mathcal{P}_{A_2}|$. Substituting $wt(\mathbf{c}_1) = 1 + |\mathcal{P}_{A_1}|$ and $wt(\mathbf{c}_2) = 1 + |\mathcal{P}_{A_2}|$ gives:
\[
wt(\mathbf{c}_1 + \mathbf{c}_2) = (1 + |\mathcal{P}_{A_1}|) + (1 + |\mathcal{P}_{A_2}|) - 2(1 + |\mathcal{P}_{A_1} \cap \mathcal{P}_{A_2}|)
\]
\[
= |\mathcal{P}_{A_1}| + |\mathcal{P}_{A_2}| - 2|\mathcal{P}_{A_1} \cap \mathcal{P}_{A_2}|.
\]
Solving for $|\mathcal{P}_{A_1} \cap \mathcal{P}_{A_2}|$ yields the exact identity. Substituting $wt(\mathbf{c}_1), wt(\mathbf{c}_2) \le w_{\max}$ and $wt(\mathbf{c}_1 + \mathbf{c}_2) \ge w_{\min} = \theta(2^m-1)$ establishes the upper bound.
\end{proof}

Finally, to demonstrate these mathematical results on a concrete construction, we evaluate the complete access distribution over a 8-dimensional field extension.

\begin{example}
Let $n = 8$ ($m = 4$) and $N = 2^8 - 1 = 255$. Let $\mathcal{M}_{(\psi_a)}$ be the $[255, 12, 106]$ minimal binary linear code constructed over a participant set of size $|\mathcal{P}| = 254$.

By Theorem~\ref{Thm:SSS_Construction}, the dual paths induce a secret sharing scheme with exactly $|\Gamma| = 2^{8+3} = 2048$ minimal access sets. The dual weight spectrum of $\mathcal{M}_{(\psi_a)}$ yields the explicit set-size distribution profile for minimal access sets via $|\mathcal{P}_A| = w - 1$. For $n = 8$, the operational span of minimal access set cardinalities satisfies:
\[
\Delta = |\mathcal{P}_{A,\max}| - |\mathcal{P}_{A,\min}| \ge 2^{8-1} + 2^{4-2} - 1 = 128 + 4 - 1 = 131.
\]
Furthermore, by Proposition~\ref{prop:information_rate}, the code rate is $R = \frac{12}{255} \approx 0.0471$. Compared to previous $(n+2)$-dimensional constructions, this scheme achieves a relative information rate enhancement of:
\[
\frac{2}{n+2} = \frac{2}{8+2} = \frac{2}{10} = 20\%.
\]
\end{example}



\section{Conclusion}
In this paper, we constructed a new family of $[2^n-1, n+4]$ minimal binary linear codes using four-component Boolean functions supported on partial spreads. By evaluating character sums across orthogonal subspace intersections, we derived their explicit weight distributions and established necessary and sufficient minimality criteria, proving that the family yields minimal codes that strictly violate the classical Ashikhmin--Barg ceiling. When applied to Massey's secret sharing framework, these codes induce ideal multi-threshold access structures that overcome lower-dimensional limits ($K \le n+3$) by providing a quadrupled authorization density ($|\Gamma| = 2^{n+3}$), a broad operational span ($\Delta \ge 2^{n-1} + 2^{m-2} - 1$), information-theoretic coalition immunity ($P_{\text{cheat}} = 1/2$), bounded participant overlaps, and a $20\%$ throughput enhancement over existing $(n+2)$-dimensional schemes.

Looking ahead, several promising research directions emerge from this framework. Key extensions include generalizing this four-component construction to non-binary finite fields $\mathbb{F}_q$, investigating subfield codes derived from higher-dimensional partial spreads, and designing multi-secret sharing protocols and secure multi-party computation primitives that leverage the multi-threshold properties of these dual access structures.

\begin{thebibliography}{99}

\bibitem{AshikhminBarg1998}
A.~Ashikhmin and A.~Barg, \textit{Minimal vectors in linear codes}, IEEE Trans. Inf. Theory 44(5), 2010--2017, 1998.

\bibitem{CarletDingYuan2005}
C.~Carlet, C.~Ding, and J.~Yuan, \textit{Linear codes from perfect nonlinear mappings and their secret sharing schemes}, IEEE Trans. Inf. Theory 51(6), 2089--2102, 2005.

\bibitem{Ding2016}
C.~Ding, \textit{A construction of binary linear codes from Boolean functions}, Discrete Math. 339(9), 2288--2303, 2016.

\bibitem{DingHengZhou2018}
C.~Ding, Z.~Heng, and Z.~Zhou, \textit{Minimal binary linear codes}, IEEE Trans. Inf. Theory 64(10), 6536--6545, 2018.

\bibitem{HengDingZhou2018}
Z.~Heng, C.~Ding, and Z.~Zhou, \textit{Minimal linear codes over finite fields}, Finite Fields Their Appl. 54, 176--196, 2018.

\bibitem{LiYue2020}
X.~Li and Q.~Yue, \textit{Four classes of minimal binary linear codes with $w_{\min}/w_{\max}<\frac12$ derived from Boolean functions}, Des. Codes Cryptogr. 88(2), 257--271, 2020.

\bibitem{LingXing2004}
S.~Ling and C.~Xing, \textit{Coding Theory: A First Course}, Cambridge University Press, Cambridge, 2004.

\bibitem{LiuLiao2022}
H.~Liu and Q.~Liao, \textit{A New Constructions of Minimal Binary Linear Codes}, 2022. \url{https://arxiv.org/pdf/2201.02981}

\bibitem{LuWuCao2021}
W.~Lu, X.~Wu, and X.~Cao, \textit{The parameters of minimal linear codes}, Finite Fields Their Appl. 71, 101799, 2021.

\bibitem{Massey1993}
J.~L.~Massey, \textit{Minimal codewords and secret sharing}, in Proc. 6th Joint Swedish-Russian Int. Workshop on Inf. Theory, M\"olle, Sweden, pp. 276--279, 1993.

\bibitem{MesnagerQianCaoYuan2023}
S.~Mesnager, L.~Qian, X.~Cao, and M.~Yuan, \textit{Several families of binary minimal linear codes from two-to-one functions}, IEEE Trans. Inf. Theory 69(5), 3285--3301, 2023.

\bibitem{YuanDing2005}
J.~Yuan and C.~Ding, \textit{Secret sharing schemes from three classes of linear codes}, IEEE Trans. Inf. Theory 52(1), 206--212, 2005.

\end{thebibliography}

\end{document}
